% chapter6.tex
% ============================================================
%  فصل ۶: راهنمای اجرایی برای ایران
% ============================================================

\chapter{نقشه‌ی راه: طراحی نهادی برای ایران}

\begin{keybox}[خلاصه‌ی فصل]
این فصل، بر اساس مبانی نظری و آموزه‌های تطبیقی فصول پیشین، یک \keyword{چارچوب اجرایی} برای
مدیریت تنوع قومی–ملّی در ایران ارائه می‌دهد. تأکید بر \concept{واقع‌گرایی}، \concept{تدریجی‌بودن}
و \concept{جامعیت} است.
\end{keybox}

% ============================================================
\section{اصول راهنما}
% ============================================================

\begin{goldbox}[هشت اصل بنیادین]
\begin{enumerate}[label=\textcolor{AccentGold}{\bfseries\Roman*.}]
  \item \keyword{یکپارچگی سرزمینی:}
    هیچ طرحی نباید تمامیت ارضی ایران را به خطر اندازد.
  \item \keyword{دموکراسی و حقوق بشر:}
    مبنای هر اصلاحی باید حقوق فردی و جمعیِ شهروندان باشد.
  \item \keyword{به‌رسمیت‌شناسی تنوع:}
    تنوع قومی، زبانی و فرهنگی باید به‌رسمیت شناخته و از آن حمایت شود.
  \item \keyword{تدریجی‌بودن:}
    اصلاحات باید مرحله‌به‌مرحله و با ارزیابی مستمر اجرا شوند.
  \item \keyword{تقارن / عدم تقارن هوشمند:}
    واحدهای مختلف می‌توانند اختیارات متفاوت داشته باشند (فدرالیسم نامتقارن).
  \item \keyword{عدالت توزیعی:}
    توسعه‌ی اقتصادی عادلانه پیش‌شرط هر اصلاح سیاسی است.
  \item \keyword{مشارکت و اجماع:}
    طراحی نهادی باید حاصل گفت‌وگوی ملّی و نه تحمیل از بالا باشد.
  \item \keyword{سکولاریسم شمول‌گرا:}
    جدایی دین از حکومت ضمن احترام به باورهای دینی همه‌ی گروه‌ها.
\end{enumerate}
\end{goldbox}

% ============================================================
\section{مرحله‌ی اول: اصلاحات فوری (کوتاه‌مدت: ۱–۳ سال)}
% ============================================================

\begin{figure}[htbp]
\centering
\begin{tikzpicture}[
  node distance=0.6cm,
  action/.style={
    draw=PrimaryMid, fill=BgLightBlue, rounded corners=6pt,
    minimum width=12cm, minimum height=1.2cm, align=right,
    font=\small, text=TextDark,
    text width=11.5cm
  },
  phase/.style={
    fill=#1, text=white, rounded corners=4pt,
    minimum width=13cm, minimum height=1cm,
    font=\bfseries\large, align=center
  }
]
  \node[phase=WarningOrange] (ph1) {مرحله‌ی اول: اصلاحات فوری (۱–۳ سال)};
  \node[action, below=0.4cm of ph1] (a1)
    {\textcolor{PrimaryDark}{\bfseries ۱.۱.}
     اجرای کامل اصل ۱۵ قانون اساسی: آموزش رسمی به زبان‌های محلی در مدارس مناطق اقلیت‌نشین};
  \node[action, below=0.15cm of a1] (a2)
    {\textcolor{PrimaryDark}{\bfseries ۱.۲.}
     تأسیس شبکه‌های رادیو و تلویزیون استانی به زبان‌های محلی با بودجه‌ی کافی};
  \node[action, below=0.15cm of a2] (a3)
    {\textcolor{PrimaryDark}{\bfseries ۱.۳.}
     تصویب قانون ضد تبعیض قومی و زبانی با مکانیسم شکایت و پیگرد};
  \node[action, below=0.15cm of a3] (a4)
    {\textcolor{PrimaryDark}{\bfseries ۱.۴.}
     آغاز برنامه‌ی ویژه‌ی توسعه‌ی اقتصادی مناطق محروم (سیستان و بلوچستان، کردستان، خوزستان)};
  \node[action, below=0.15cm of a4] (a5)
    {\textcolor{PrimaryDark}{\bfseries ۱.۵.}
     افزایش اختیارات شوراهای شهر و استان (بازنگری قانون شوراها)};
  \node[action, below=0.15cm of a5] (a6)
    {\textcolor{PrimaryDark}{\bfseries ۱.۶.}
     تشکیل «کمیسیون ملّی گفت‌وگوی قومی» با مشارکت نمایندگان همه‌ی گروه‌ها};
\end{tikzpicture}
\caption{اقدامات مرحله‌ی اول (کوتاه‌مدت)}
\label{fig:phase1}
\end{figure}

% ============================================================
\section{مرحله‌ی دوم: تمرکززدایی ساختاری (میان‌مدت: ۳–۷ سال)}
% ============================================================

\begin{figure}[htbp]
\centering
\begin{tikzpicture}[
  node distance=0.6cm,
  action/.style={
    draw=AccentGold, fill=BgLightGold, rounded corners=6pt,
    minimum width=12cm, minimum height=1.2cm, align=right,
    font=\small, text=TextDark,
    text width=11.5cm
  },
  phase/.style={
    fill=#1, text=white, rounded corners=4pt,
    minimum width=13cm, minimum height=1cm,
    font=\bfseries\large, align=center
  }
]
  \node[phase=AccentGold] (ph2) {مرحله‌ی دوم: تمرکززدایی ساختاری (۳–۷ سال)};
  \node[action, below=0.4cm of ph2] (b1)
    {\textcolor{PrimaryDark}{\bfseries ۲.۱.}
     ایجاد مجالس قانون‌گذاری منطقه‌ای با اختیارات مشخص
     (آموزش، بهداشت، فرهنگ، زیرساخت محلی)};
  \node[action, below=0.15cm of b1] (b2)
    {\textcolor{PrimaryDark}{\bfseries ۲.۲.}
     انتخابی‌شدن استانداران و فرمانداران (به‌جای انتصابی)};
  \node[action, below=0.15cm of b2] (b3)
    {\textcolor{PrimaryDark}{\bfseries ۲.۳.}
     تخصیص سهم مشخصی از درآمدهای ملّی (نفت، مالیات) به مناطق بر اساس فرمول عادلانه};
  \node[action, below=0.15cm of b3] (b4)
    {\textcolor{PrimaryDark}{\bfseries ۲.۴.}
     ایجاد نظام مالیاتی دوسطحی (ملّی + منطقه‌ای)};
  \node[action, below=0.15cm of b4] (b5)
    {\textcolor{PrimaryDark}{\bfseries ۲.۵.}
     تأسیس مجلس دوم (مجلس مناطق/اقوام) با نمایندگی تضمین‌شده‌ی همه‌ی گروه‌ها};
  \node[action, below=0.15cm of b5] (b6)
    {\textcolor{PrimaryDark}{\bfseries ۲.۶.}
     ایجاد دادگاه قانون اساسی مستقل برای حل اختلافات مرکز–منطقه};
\end{tikzpicture}
\caption{اقدامات مرحله‌ی دوم (میان‌مدت)}
\label{fig:phase2}
\end{figure}

% ============================================================
\section{مرحله‌ی سوم: نهادسازی فدرال (بلندمدت: ۷–۱۵ سال)}
% ============================================================

\begin{figure}[htbp]
\centering
\begin{tikzpicture}[
  node distance=0.6cm,
  action/.style={
    draw=SuccessGreen, fill=green!4, rounded corners=6pt,
    minimum width=12cm, minimum height=1.2cm, align=right,
    font=\small, text=TextDark,
    text width=11.5cm
  },
  phase/.style={
    fill=#1, text=white, rounded corners=4pt,
    minimum width=13cm, minimum height=1cm,
    font=\bfseries\large, align=center
  }
]
  \node[phase=SuccessGreen] (ph3) {مرحله‌ی سوم: نهادسازی فدرال (۷–۱۵ سال)};
  \node[action, below=0.4cm of ph3] (c1)
    {\textcolor{PrimaryDark}{\bfseries ۳.۱.}
     تدوین قانون اساسی جدید با ماهیت فدرال از طریق مجلس مؤسسان منتخب};
  \node[action, below=0.15cm of c1] (c2)
    {\textcolor{PrimaryDark}{\bfseries ۳.۲.}
     تعیین واحدهای فدرال بر اساس ترکیب جغرافیایی–فرهنگی (نه صرفاً قومی)};
  \node[action, below=0.15cm of c2] (c3)
    {\textcolor{PrimaryDark}{\bfseries ۳.۳.}
     تنظیم فهرست‌های سه‌گانه‌ی صلاحیت: فدرال، منطقه‌ای، مشترک};
  \node[action, below=0.15cm of c3] (c4)
    {\textcolor{PrimaryDark}{\bfseries ۳.۴.}
     تضمین حقوق اساسی شهروندی فراقومی در سطح فدرال};
  \node[action, below=0.15cm of c4] (c5)
    {\textcolor{PrimaryDark}{\bfseries ۳.۵.}
     ایجاد مکانیسم‌های فدرالیسم مالی عادلانه (صندوق انتقالی، بودجه‌ی تعدیلی)};
  \node[action, below=0.15cm of c5] (c6)
    {\textcolor{PrimaryDark}{\bfseries ۳.۶.}
     همه‌پرسی ملّی برای تصویب قانون اساسی جدید};
\end{tikzpicture}
\caption{اقدامات مرحله‌ی سوم (بلندمدت)}
\label{fig:phase3}
\end{figure}

% ============================================================
\section{ساختار نهادی پیشنهادی}
% ============================================================

\begin{figure}[htbp]
\centering
\begin{tikzpicture}[
  scale=0.85, transform shape,
  inst/.style={
    draw=PrimaryMid, fill=BgLightBlue, rounded corners=8pt,
    minimum width=4cm, minimum height=1.8cm, align=center,
    font=\small\bfseries, text=PrimaryDark,
    drop shadow={opacity=0.12}
  },
  sub/.style={
    draw=AccentGold, fill=BgLightGold, rounded corners=6pt,
    minimum width=3.2cm, minimum height=1.2cm, align=center,
    font=\footnotesize, text=PrimaryDark
  },
  arr/.style={-{Stealth[length=5pt]}, thick, PrimaryMid},
  darr/.style={-{Stealth[length=5pt]}, thick, AccentGold, dashed}
]
  % Federal Level
  \node[fill=PrimaryDark, text=white, rounded corners=4pt,
        minimum width=14cm, minimum height=1cm, font=\large\bfseries]
    (flabel) at (0,8) {سطح فدرال};

  \node[inst] (president) at (-5,6) {رئیس‌جمهور\\(منتخب مستقیم)};
  \node[inst] (parl1) at (0,6) {مجلس ملّی\\(نمایندگان مردم)};
  \node[inst] (parl2) at (5,6) {مجلس مناطق\\(نمایندگان واحدها)};
  \node[inst] (court) at (0,3.8) {دادگاه قانون\\اساسی};
  \node[inst] (cabinet) at (-5,3.8) {هیئت دولت\\فدرال};

  % Regional Level
  \node[fill=AccentGold, text=PrimaryDark, rounded corners=4pt,
        minimum width=14cm, minimum height=1cm, font=\large\bfseries]
    (rlabel) at (0,1.8) {سطح منطقه‌ای (واحدهای فدرال)};

  \node[sub] (rgov) at (-4,0) {فرماندار منتخب\\منطقه};
  \node[sub] (rparl) at (0,0) {مجلس قانون‌گذاری\\منطقه‌ای};
  \node[sub] (rcourt) at (4,0) {دادگستری\\منطقه‌ای};

  % Local Level
  \node[fill=SuccessGreen, text=white, rounded corners=4pt,
        minimum width=14cm, minimum height=1cm, font=\large\bfseries]
    (llabel) at (0,-1.5) {سطح محلی};

  \node[sub, draw=SuccessGreen, fill=green!5] (city) at (-3,-3)
    {شهرداری‌ها\\و شوراهای شهر};
  \node[sub, draw=SuccessGreen, fill=green!5] (village) at (3,-3)
    {دهیاری‌ها\\و شوراهای روستا};

  % Arrows
  \draw[arr] (president) -- (cabinet);
  \draw[arr] (parl1) -- (court);
  \draw[arr] (parl2) -- (court);
  \draw[darr] (cabinet) -- (rgov);
  \draw[darr] (parl1.south) -- ++(0,-0.5) -| (rparl);
  \draw[darr] (court) -- (rcourt);
  \draw[SuccessGreen, thick, dashed, -{Stealth[length=4pt]}]
    (rgov) -- (city);
  \draw[SuccessGreen, thick, dashed, -{Stealth[length=4pt]}]
    (rparl) -- (village);

  % Bidirectional for Senate
  \draw[arr, PrimaryDark] (parl2.south east) -- ++(0.5,-0.5) -- (rcourt.north east)
    node[midway, right, font=\scriptsize\color{PrimaryDark}] {نمایندگی};
\end{tikzpicture}
\caption{ساختار نهادی پیشنهادی سه‌سطحی برای ایران فدرال}
\label{fig:proposed-structure}
\end{figure}

% ============================================================
\section{تقسیم صلاحیت‌ها}
% ============================================================

\begin{sidewaystable}
\centering
\caption{پیشنهاد تقسیم صلاحیت‌ها در نظام فدرال ایران}
\label{tab:competence-division}
\small
\rowcolors{2}{BgLightBlue!30}{white}
\begin{tabularx}{\linewidth}{
  >{\raggedright\arraybackslash}X
  >{\raggedright\arraybackslash}X
  >{\raggedright\arraybackslash}X
}
\toprule
\rowcolor{PrimaryDark}
\textcolor{white}{\bfseries صلاحیت فدرال (انحصاری)} &
\textcolor{white}{\bfseries صلاحیت مشترک} &
\textcolor{white}{\bfseries صلاحیت منطقه‌ای (انحصاری)} \\
\midrule
دفاع ملّی و ارتش &
بهداشت عمومی &
آموزش (محتوای منطقه‌ای + زبان) \\
سیاست خارجی و دیپلماسی &
محیط زیست &
فرهنگ و میراث فرهنگی محلی \\
پول و بانکداری مرکزی &
انرژی (تولید و توزیع) &
گردشگری منطقه‌ای \\
گمرک و تجارت خارجی &
حمل‌ونقل بین‌منطقه‌ای &
حمل‌ونقل درون‌منطقه‌ای \\
حقوق شهروندی و قانون مدنی پایه &
آموزش عالی &
کشاورزی و منابع طبیعی محلی \\
دادگاه قانون اساسی &
تأمین اجتماعی &
مسکن و شهرسازی \\
مالیات‌های فدرال &
صنعت و معدن &
مالیات‌های منطقه‌ای \\
مهاجرت و تابعیت &
امنیت داخلی (هماهنگی) &
پلیس محلی \\
\bottomrule
\end{tabularx}
\end{sidewaystable}

% ============================================================
\section{فدرالیسم مالی: مدل پیشنهادی}
% ============================================================

\begin{figure}[htbp]
\centering
\begin{tikzpicture}[
  node distance=1cm,
  fund/.style={
    draw=PrimaryMid, fill=BgLightBlue, rounded corners=8pt,
    minimum width=4cm, minimum height=2cm, align=center,
    font=\small\bfseries, text=PrimaryDark
  },
  flow/.style={-{Stealth[length=6pt]}, thick, #1}
]
  % Revenue sources
  \node[fund] (natrev) at (0,5) {درآمدهای ملّی\\(نفت، گمرک،\\مالیات فدرال)};

  % Distribution
  \node[fund, fill=BgLightGold, draw=AccentGold] (pool) at (0,2)
    {صندوق توزیع\\ملّی};

  % Destinations
  \node[fund] (fed) at (-5,-1) {بودجه‌ی\\فدرال\\(۵۰\%)};
  \node[fund] (reg) at (0,-1) {سهم مناطق\\بر اساس فرمول\\(۳۵\%)};
  \node[fund] (eq) at (5,-1) {صندوق\\تعدیل\\(۱۵\%)};

  % Formula detail
  \node[draw=AccentGold, fill=BgLightGold, rounded corners=4pt,
        minimum width=5.5cm, align=center, font=\footnotesize, text=TextDark,
        below=0.5cm of reg]
    (formula) {فرمول = f(جمعیت، وسعت،\\
    شاخص محرومیت، درآمد سرانه)};

  \node[draw=SuccessGreen, fill=green!5, rounded corners=4pt,
        minimum width=5cm, align=center, font=\footnotesize, text=TextDark,
        below=0.5cm of eq]
    (eqdetail) {تخصیص ویژه به\\
    مناطق محروم و اقلیت‌نشین\\
    (سیستان، کردستان، خوزستان)};

  \draw[flow=PrimaryMid] (natrev) -- (pool);
  \draw[flow=PrimaryDark] (pool) -- (fed);
  \draw[flow=AccentGold] (pool) -- (reg);
  \draw[flow=SuccessGreen] (pool) -- (eq);
  \draw[AccentGold, dashed] (reg) -- (formula);
  \draw[SuccessGreen, dashed] (eq) -- (eqdetail);

  % Regional own revenue
  \node[fund, draw=WarningOrange, fill=WarningOrange!8] (ownrev) at (-5,-4)
    {درآمدهای اختصاصی\\منطقه‌ای\\(مالیات محلی،\\عوارض)};

  \draw[flow=WarningOrange, dashed] (ownrev.east) -- ++(2,0)
    node[above, midway, font=\scriptsize\color{WarningOrange}] {مکمل};
\end{tikzpicture}
\caption{مدل فدرالیسم مالی پیشنهادی}
\label{fig:fiscal-federalism}
\end{figure}

% ============================================================
\section{سیاست زبانی پیشنهادی}

<!-- POSSIBLE OVERLAP DETECTED (similarity: 90%) - REVIEW BELOW -->

ادامه‌ی پروژه را از بخش سیاست زبانی فصل ششم پی می‌گیرم:

---

## ادامه‌ی `chapter6.tex`